%Packages and document class
\documentclass[epj]{svjour}
\usepackage{textcomp}
\usepackage{amssymb,amsmath,verbatim,mathtools,needspace,enumitem,etoolbox,graphicx,microtype,afterpage,bigints,gensymb,tabularx,xspace,siunitx}
\usepackage{aasmacros}
\usepackage{widetext}
\usepackage{graphicx}% Include figure files
\usepackage{dcolumn}% Align table columns on decimal point
\usepackage{bm}% bold math
\usepackage{lipsum}
\usepackage[usenames, dvipsnames, table]{xcolor}
\usepackage{physics}
\usepackage{soul}
%\usepackage{natbib}
\usepackage{graphics}
\definecolor{linkcolor}{rgb}{0.0,0.3,0.5}
\usepackage[colorlinks = true,
            linkcolor = linkcolor,
            urlcolor  = linkcolor,
            citecolor = linkcolor,
            anchorcolor = linkcolor]{hyperref}
\usepackage{array}
\usepackage{orcidlink}


% Comment commands
\newcommand{\rb}[1]{\textcolor{red}{#1}}
\newcommand{\citeme}{\textcolor{blue}{CIT.}}
\newcommand{\citethis}[1]{\textcolor{blue}{CIT: #1}}
\newcommand{\ssim}{\mathchar"5218\relax\,}
\newcommand{\ttimes}{\!\times\!\relax}

% Math commands
\newcommand{\sub}[1]{_{\text{#1}}}
\newcommand{\super}[1]{^{\text{#1}}}
\newcommand{\uvec}[1]{\bm{\hat{#1}}}
\newcommand{\dvec}[1]{\dot{\bm{#1}}}
\newcommand{\ddvec}[1]{\ddot{\bm{#1}}}
\newcommand{\dddvec}[1]{\dddot{\bm{#1}}}
\newcommand{\duvec}[1]{\dot{\bm{\hat{#1}}}}
\newcommand{\dduvec}[1]{\ddot{\bm{\hat{#1}}}}
\newcommand{\ddduvec}[1]{\dddot{\bm{\hat{#1}}}}
\newcommand{\ord}[1]{\mathcal{O} \left( #1 \right)}

%affiliation commands
\newcommand{\bicocca}{Dipartimento di Fisica “G. Occhialini”, Università degli Studi di Milano-Bicocca, Piazza della Scienza 3, 20126 Milano, Italy}
\newcommand{\infn}{INFN, Sezione di Milano-Bicocca, Piazza della Scienza 3, 20126 Milano, Italy}
\newcommand{\bham}{Institute for Gravitational Wave Astronomy \& School of Physics and Astronomy, University of Birmingham, Birmingham, B15 2TT, UK}
\newcommand{\mpa}{Max-Planck-Institut f{\"u}r Astrophysik, Karl-Schwarzschild-Straße 1, 85741 Garching, Germany}

% SI unit options
\AtBeginDocument{\RenewCommandCopy\qty\SI}
\ExplSyntaxOn
\msg_redirect_name:nnn { siunitx } { physics-pkg } { none }
\ExplSyntaxOff
% Make the ranges with siunitx shown with a dash an unit only shown once
% % astrophysical units
\DeclareSIUnit\solarmass{\ensuremath{M_{\odot}}}
\DeclareSIUnit\parsec{pc}
\DeclareSIUnit\lightspeed{$c$}
\DeclareSIUnit\year{yr}
% redeclared units
\DeclareSIUnit\arcsecond{as}
\DeclareSIUnit\astronomicalunit{AU}
\DeclareSIUnit\clight{\ensuremath c}
%

% If you need just the number to write the rest, use 
%\labelcref{your-label}
%
\date{Received: date / Revised version: date}

\begin{document}
\sisetup{range-phrase=-, range-units=single}

\title{Title}

\author{Riccardo Buscicchio\inst{1,2,4}\thanks{\href{mailto:riccardo.buscicchio@unimib.it}{riccardo.buscicchio@unimib.it}}%~\orcidlink{0000-0002-7387-6754}%
\and 
John Doe\inst{4}%~\orcidlink{0000-0001-5438-9152}
}                     % Do not remove

\institute{\bicocca \and \infn \and \mpa \and \bham}

%Shorter header 
\titlerunning{Header title}
\authorrunning{Buscicchio et al.}

%
\date{Received: date / Revised version: date}
% The correct dates will be entered by Springer
%

\abstract{
An abstract
%
\PACS{
      {04.80.Nn, 95.55.Ym}{Gravitational wave detectors}   \and
      {95.85.Sz}{Gravitational wave astronomical observations}
     } % end of PACS codes
} %end of abstract
\maketitle

\section{\label{sec:a} Introduction}
Intro with a range using SI units $\SIrange{e-4}{e-1}{\hertz}$.
A citation Ref.~\cite{2023LRR....26....2A}. 

The paper is organized as follows.
In~\ref{sec:a}, we do such and such.
Another type of citation~\cite{2023JCAP...08..034B}. 
Some dimensionful quantities \SI{3}{\kilo\parsec}. 
Some other \SI{13.5}{\giga\year}.
Another quantity with scale factor \SI{8.2e10}{\solarmass}. 
A pure number \num{26e6}. 

\begin{equation}
    \dot{f} = \frac{96 \pi^{8/3}}{5 c^5}  \left( G {\cal M}_c \right)^{5/3} f^{11/3}, \label{eq:fdot}
\end{equation}
A percentage \SI{1}{\percent} and some dimensionful ranges $f>\SIrange{2}{3}{\milli\hertz}$ or $f<\SIrange{2}{3}{\milli\hertz}$.

\section{\label{sec:b} Methodology}

\begin{figure}[t]
    \centering
    \includegraphics[width=1.0\columnwidth]{example-image-a}
    \caption{A single column figure}\label{fig:1}
\end{figure}

\clearpage
\bibliographystyle{epjc} %apsrev4-2
\bibliography{main}% Produces the bibliography via BibTeX.

\end{document}

% end of file template.tex

